% =====================================================================
\chapter{Flujo residual y posición de la normalización}
\label{tema:2}
% =====================================================================

En el módulo~11 se estudiaron las conexiones residuales y la normalización para una red genérica. Este tema las aplica al bloque del \cref{tema:1}: desarrolla el flujo residual de un Transformer como suma de $2L$ aportaciones, compara las variantes pre-norm y post-norm y justifica la normalización final y la inicialización escalada. Se sigue~\cite[§7.2]{jm3}, la discusión de la práctica~1 de CS336~\cite[§3.4]{cs336}, el análisis de Xiong \emph{et al.}~\cite{xiong}, la §2.3 de~\cite{gpt2} y la interpretación de Elhage \emph{et al.}~\cite{elhage}.

\section{El flujo residual}

Para simplificar la notación, se numeran las $2L$ subcapas del Transformer en el orden en que se aplican: la subcapa $k=2\ell-1$ es la atención del bloque $\ell$ y la subcapa $k=2\ell$ es su MLP. Sea $F_k$ la función que calcula la subcapa $k$ \emph{incluida su normalización de entrada}: $F_{2\ell-1}=\MHA\circ\Norm_1$ y $F_{2\ell}=\MLP\circ\Norm_2$.

\begin{definicion}[Flujo residual]\label{def:flujo}
El \emph{flujo residual} es la sucesión de matrices $R^{0}=H^{0}$ y $R^{k}=R^{k-1}+F_k(R^{k-1})$, $k=1,\dots,2L$. Se cumple $R^{2\ell}=H^{\ell}$ y $R^{2\ell-1}=A^{\ell}$, la matriz intermedia del bloque $\ell$. La fila $t$ de $R^{k}$ es el vector del flujo del token $t$.
\end{definicion}

\begin{proposicion}[Desarrollo telescópico]\label{prop:telescopica}
En un Transformer pre-norm,
\[
H^{L}=H^{0}+\sum_{k=1}^{2L}F_k\bigl(R^{k-1}\bigr).
\]
\end{proposicion}

La \cref{prop:telescopica} sustenta la lectura del flujo como un \emph{canal de comunicación}~\cite{elhage}: ninguna subcapa sobrescribe el flujo; cada una lee de él (tras normalizar) y le \emph{suma} su salida. De ahí se siguen tres consecuencias.
\begin{enumerate}
  \item Una subcapa puede escribir en una dirección de $\R^{d}$ que lee otra subcapa muchas capas después. Esta descomposición es la base de la interpretabilidad mecanicista (módulo~23).
  \item Toda subcapa debe producir vectores de $\R^{d}$, pues su salida se suma al flujo. Esa es la función de $W_O$ en la atención y de $\Wdown$ en el MLP.
  \item Si una subcapa aprende a escribir casi cero, el modelo se comporta como uno con una subcapa menos.
\end{enumerate}

\addtocounter{definicion}{1}%

\begin{proposicion}[Jacobiano del flujo]\label{prop:jac-flujo}
Sea $J_k$ el jacobiano de $F_k$ evaluado en $R^{k-1}$, con todas las matrices vectorizadas. Entonces
\[
\frac{\partial H^{L}}{\partial H^{0}}=\prod_{k=2L}^{1}\bigl(I+J_k\bigr)=I+\sum_{k=1}^{2L}J_k+\sum_{k>k'}J_kJ_{k'}+\cdots+J_{2L}\cdots J_1,
\]
donde el producto se ordena de izquierda a derecha con $k$ decreciente.
\end{proposicion}

El término $I$ del desarrollo es el \emph{camino identidad}: el gradiente de la pérdida respecto de $H^{L}$ llega a $H^{0}$ sumado, sin atravesar ninguna normalización ni producto de jacobianos. Es la versión para el Transformer del teorema del jacobiano residual del módulo~11.

\section{Pre-norm frente a post-norm}

\begin{definicion}[Post-norm]\label{def:postnorm}
El \emph{bloque post-norm} del Transformer original~\cite[§3.1]{vaswani} aplica la normalización después de la suma residual:
\[
A=\Norm_1\bigl(H+\MHA(H)\bigr),\qquad H'=\Norm_2\bigl(A+\MLP(A)\bigr).
\]
\end{definicion}

\begin{proposicion}[No hay camino identidad en post-norm]\label{prop:postnorm}
En un Transformer post-norm con normalización RMSNorm sin ganancia, $\Norm(x)=x/\RMS(x)$, se cumple $\RMS(h^{\ell}_t)=1$ para toda fila $t$ y toda capa $\ell\ge 1$, y el jacobiano de la salida respecto de la entrada es
\[
\frac{\partial H^{L}}{\partial H^{0}}=\prod_{k=2L}^{1}N_k\,(I+J_k),
\]
donde $N_k$ es el jacobiano de la normalización $k$. Cada $N_k$ es singular.
\end{proposicion}

En post-norm, toda la señal, y todo el gradiente, atraviesa $2L$ normalizaciones; no existe un término $I$ aislado como en la \cref{prop:jac-flujo}. Xiong \emph{et al.}~\cite[§3.3, teorema~1]{xiong} cuantifican la diferencia en la inicialización: bajo hipótesis de concentración de las normas, la norma del gradiente de los pesos de la última capa es del orden de $d\sqrt{\ln d}$ en post-norm, sin dependencia de $L$, y del orden de $d\sqrt{\ln d/L}$ en pre-norm. Gradientes grandes en las capas altas obligan a un calentamiento (\emph{warmup}) largo de la tasa de aprendizaje en post-norm; pre-norm se entrena de forma estable con un calentamiento corto o nulo.

\begin{table}[ht]
\centering\small
\caption{Comparación de las dos posiciones de la normalización.}\label{tab:prepost}
\begin{tabular}{@{}lll@{}}
\toprule
Aspecto & Post-norm (2017) & Pre-norm (GPT-2, LLaMA) \\
\midrule
Ecuación & $\Norm(x+F(x))$ & $x+F(\Norm(x))$ \\
Camino identidad & no (\cref{prop:postnorm}) & sí (\cref{prop:jac-flujo}) \\
Norma del flujo & fija & crece con la profundidad \\
Calentamiento & imprescindible & tolerante \\
Normalización final & innecesaria & necesaria \\
\bottomrule
\end{tabular}
\end{table}

\section{El precio de pre-norm: crecimiento del flujo}

\begin{proposicion}[Crecimiento del flujo en un Transformer]\label{prop:crecimiento}
Supóngase que, en la inicialización, las aportaciones $F_k(R^{k-1})_t$ de una fila son vectores aleatorios de media nula, incorrelados entre sí y con $h^{0}_t$, con $\E\norm{F_k(R^{k-1})_t}^{2}=s^{2}$. Entonces
\[
\E\norm{h^{L}_t}^{2}=\E\norm{h^{0}_t}^{2}+2L\,s^{2}.
\]
\end{proposicion}

La \cref{prop:crecimiento} tiene dos consecuencias prácticas. Primera: como la norma del flujo crece, el vector $h^{L}_t$ debe normalizarse antes del unembedding; de lo contrario, la escala de los logits dependería de la profundidad. Es la normalización final $\Norm_f$ de la \cref{def:transformer}, que GPT-2 introdujo junto con la posición pre-norm~\cite[§2.3]{gpt2} y que Jurafsky y Martin señalan como la capa adicional que exige pre-norm~\cite[§7.2.3]{jm3}. Segunda: GPT-2 multiplica en la inicialización los pesos de las capas que escriben en el flujo ($W_O$ y la última matriz del MLP) por $1/\sqrt{N}$, con $N=2L$ el número de capas residuales; así cada $s^{2}$ se divide por $2L$ y la norma esperada del flujo final queda acotada independientemente de la profundidad.

\addtocounter{definicion}{1}%

\addtocounter{definicion}{1}%

